\section{Introduction and Motivation}
\label{sec:intro}

Self-attention \citep{vaswani2017attention} compares every token with every earlier token. This gives the
Transformer an exact, content-addressable memory of its entire context, but at a price: training costs
$O(T^2)$ time in the sequence length $T$, and generation must keep a key--value (KV) cache whose size
grows linearly with the context. For long contexts the KV cache, not the arithmetic, becomes the
bottleneck of deployment.

A family of \emph{linear recurrent} sequence models removes this cost. State-space models (SSMs) such as
Mamba \citep{gu2023mamba} and Mamba-2 \citep{dao2024ssd}, and linear-attention models such as Gated
DeltaNet \citep{yang2025gated}, compress the whole past into a state of fixed size, updated once per token.
Training runs in $O(T)$ with a parallel scan, and generation needs constant memory per sequence. These
layers are no longer a research curiosity: recent production models are hybrids that replace most of
their attention layers with them, for example Kimi Linear \citep{kimi2025linear}, which interleaves three
delta-rule layers with every attention layer.

The price of a fixed-size state is that the memory must \emph{fade}: information is compressed, decayed
and overwritten. Which pattern-recognition abilities survive this compression, and which do not, is the
question of this project. Two abilities are particularly revealing.
\begin{itemize}
  \item \textbf{Associative recall}: storing many key--value facts and retrieving one by its key. Attention
        does this by direct lookup; a recurrent model must fit all facts into its state.
  \item \textbf{State tracking}: maintaining a running summary that changes with every input, such as the
        parity of a bit string or the composition of a sequence of permutations. Here the question is not
        how much the state can hold, but which updates the recurrence can express at all.
\end{itemize}
Both are pattern-recognition problems in the classical sense: a sequence must be mapped to a label, and
the model must discover a rule from examples rather than memorise them. We probe them with controlled
synthetic tasks, which isolate a single ability and allow testing on sequences longer than any seen in
training, the cleanest check that a rule rather than a shortcut was learned.

\paragraph{Research questions.}
\begin{enumerate}
  \item[\textbf{RQ1}] How much associative recall can a fixed-size state support, and how does this depend
        on the size of the state and on the update rule?
  \item[\textbf{RQ2}] Which state transitions can track state? We test parity ($\mathbb{Z}_2$), counting
        modulo 3 ($\mathbb{Z}_3$) and the non-abelian permutation group $S_3$, the negative-eigenvalue fix
        of \citet{grazzi2025unlocking}, and two recent extensions: complex-valued (rotational) transitions
        in the spirit of Mamba-3 \citep{mamba3} and products of Householder matrices
        (DeltaProduct, \citealp{siems2025deltaproduct}).
  \item[\textbf{RQ3}] What does the linear-time recurrence buy in latency and memory on a consumer GPU, and
        how much does the quality of the GPU kernel matter?
\end{enumerate}

\paragraph{Contributions.}
\begin{itemize}
  \item From-scratch implementations of a selective SSM, Gated DeltaNet, a rotational SSM, DeltaProduct
        and a Transformer baseline. Each recurrent mixer has a token-by-token reference and a
        chunkwise-parallel training algorithm, verified against each other to float64 precision, and
        constant-memory decoding verified against the parallel forward pass.
  \item A Triton GPU kernel for the SSM scan and an efficiency study on an RTX~4050 laptop GPU.
  \item A controlled comparison on state tracking (8 models $\times$ 3 groups $\times$ 3 seeds) and
        associative recall, all on one 6\,GB GPU.
  \item An analysis of four failure cases, including a direct inspection of learned weights showing that a
        model which \emph{can} represent the solution learned a shortcut instead.
\end{itemize}

\paragraph{Scope.} All experiments use small models (0.2--0.44M parameters, two layers) on synthetic
tasks. The aim is not state-of-the-art accuracy but a clean, mechanistic comparison whose conclusions can
be traced back to the mathematics of each recurrence.
